\documentclass[12pt]{article}
\usepackage{amsmath,amssymb}
\newcommand{\Dr}{\operatorname{Dr}}
\newcommand{\Trop}{\operatorname{Trop}}
\newcommand{\rk}{\operatorname{rk}}
\newcommand{\cl}{\operatorname{cl}}
\newcommand{\cL}{\mathcal L}
\begin{document}
\section*{Research notes}

Problem: Given a matroid $M$ and the relative Dressian $\Dr(M)$ of
tropical linear subspaces contained in $\Trop(M)$, suppose
$\cL(M)$ has an order-reversing involution $F\mapsto F^\perp$.
The embedded copy of $\cL(M)^{\rm op}$ sends a flat $F$ to the quotient
matroid $M/F\oplus U_{0,F}$.  We seek an order-reversing involution
of $\Dr(M)$ extending this flat opposition.

\subsection*{Conventions and quotient criterion}

Let $\mu_M$ be the trivial rank $r=\rk M$ valuated matroid supported
on the bases of $M$.  A rank $k$ point of $\Dr(M)$ is represented by
a projective rank $k$ valuated matroid $\nu$ with
$\nu\preceq\mu_M$, where $p\preceq q$ means that $p$ is a valuated
quotient of $q$.  Inclusion of projective tropical linear spaces is
the quotient order:
\[
\Trop(\theta)\subseteq\Trop(\nu)\Longleftrightarrow \theta\preceq\nu.
\]
Valuated matroids are considered up to adding a common scalar to all
finite Plucker coordinates.

For valuated matroids $p,q$ of ranks $a\le b$, the nested incidence
criterion is
\[
\mathcal I_{a,b}(p,q;S,T):\quad
\min_{i\in T\setminus S}\{p(S\cup i)+q(T\setminus i)\}
\text{ is attained at least twice},
\]
for $|S|=a-1$, $|T|=b+1$, $S\subseteq T$.  Together with the
valuated matroid axioms and the underlying ordinary quotient condition,
this standard nested cryptomorphic form is the quotient relation; it is
equivalent to the arbitrary-index incidence-Plucker relations.
For consecutive ranks it is the three-term relation
\[
 \min\{p(Sx)+q(Syz),p(Sy)+q(Sxz),p(Sz)+q(Sxy)\}.
\]
The convention is that, if all terms are $\infty$, the minimum is
attained at every index.

Complete flags: for a complete sequence with nonempty layers and
fixed rank $0$ and rank $r$ endpoints, adjacent three-term relations
are the local $M^\natural$-concavity conditions; by Murota's
independent-set cryptomorphism and the flag interpretation of
Brandt--Eur--Zhang/Jarra--Lorscheid, they give a complete flag of
valuated matroid quotients.

\subsection*{Modularity forced by the involution}

A finite flat lattice is ranked and semimodular.  If it has an
anti-automorphism $x\mapsto x^\perp$, ranks are reversed:
$\rk(x^\perp)=r-\rk(x)$.  Applying the anti-automorphism to
semimodularity gives the reverse inequality, hence
\[
\rk(x)+\rk(y)=\rk(x\wedge y)+\rk(x\vee y)
\]
for all $x,y$.  Therefore $\cL(M)$ is modular and
\[
(x\vee y)^\perp=x^\perp\wedge y^\perp,\qquad
(x\wedge y)^\perp=x^\perp\vee y^\perp.
\]
No simplification reduction is necessary: loops lie in
$\cl(\emptyset)$ and parallel representatives receive equal values by
flat constancy.

\subsection*{Flat constancy}

If $\nu\preceq\mu_M$ has rank $k$, then finite $\nu$-bases are
independent in $M$ because the underlying matroid of $\nu$ is a
quotient of $M$.  Also $\nu(B)$ depends only on $\cl_M(B)$ for
independent $k$-sets $B$.

Adjacent-basis proof for $0<k\le r$: compare $B=A b$ and
$B'=A b'$ with $\cl(A b)=\cl(A b')$.  Choose $C$ of size $r-k$ such
that $A b C$ is a basis of $M$; then $A b' C$ is also a basis.  In
$\mathcal I_{k,r}(\nu,\mu_M;A,A b b' C)$ only the $b,b'$ terms have
finite $\mu_M$ factor, since deleting an element of $C$ leaves the
dependent subset $A b b'$.  Thus $\nu(A b)=\nu(A b')$.  The case
$k=0$ is trivial.

Thus a rank $k$ quotient gives flat coordinates
\[
\bar\nu(X)=\nu(B)\quad\text{for }\rk X=k,\ \cl(B)=X,
\]
and $\nu(A)=\bar\nu(\cl A)$ for independent $A$, while dependent
$A$ has value $\infty$.

\subsection*{One-point completion through a given quotient}

For a single quotient $\nu\preceq\mu_M$ of rank $k$, standard
valuated truncation and relative Higgs elongation insert $\nu$ into a
complete flag:
\[
 h_j(B)=
 \begin{cases}
 \min\{\nu(C): B\subseteq C,\ |C|=k\}, & j\le k,\ |B|=j,\\
 \min\{\nu(C): C\subseteq B,\ |C|=k\}, & j\ge k,\ B\in\mathcal I_j(M),\\
 \infty, & j\ge k,\ B\notin\mathcal I_j(M).
 \end{cases}
\]
This is Murota's valuated truncation/elongation or
$M^\natural$-convex convolution with the indicator of $\mathcal I(M)$.
It gives $h_0\preceq\cdots\preceq h_r=\mu_M$ up to scalar and
$h_k=\nu$.

Endpoint check for $h_r$: for an $M$-basis $B$, let
$m(B)=\min\{\nu(C):C\subseteq B,\ |C|=k\}$.  If
$B'=B-e+f$ and $C\subseteq B$ minimizes $m(B)$, then if $e\notin C$
we have $C\subseteq B'$.  If $e\in C$, apply the incidence relation
for $\nu\preceq\mu_M$ with lower set $C-e$ and upper set $B\cup f$.
The $e$-term has value $\nu(C)$, so either it is not minimal or
there is another finite term of value at most $\nu(C)$; that term
gives a $k$-subset of $B'$.  Hence $m(B')\le m(B)$, and by symmetry
$m$ is constant on the basis graph.  After scalar normalization
$h_r=\mu_M$.

Important distinction: this one-point completion is used only to prove
that the opposition transform maps the underlying set of $\Dr(M)$ to
itself.  It is not a general partial-flag completion theorem for an
arbitrary pair $\theta\preceq\nu$; full order reversal is proved
directly by incidence duality.

\subsection*{Candidate transform and elementary opposition}

For a rank $k$ quotient $\nu$ define
\[
\nu^{\perp_M}(A)=
\begin{cases}
\bar\nu((\cl A)^\perp),& A\text{ independent in }M,\ |A|=r-k,\\
\infty,&\text{otherwise}.
\end{cases}
\]
Adding a scalar to $\nu$ adds the same scalar to $\nu^{\perp_M}$.
The formula is involutive on flat coordinates.

Local elementary opposition: Let $p\preceq q\preceq\mu_M$ with
$\rk p=d$, $\rk q=d+1$.  To check the elementary three-term relation
for $q^{\perp_M},p^{\perp_M}$, fix $A$ of size $r-d-2$ and
$x,y,z\notin A$.  Let $U=\cl(A)$ and $W=\cl(Axyz)$, with
$\delta=\rk W-\rk U$.

If $A$ is dependent or $\delta\le1$, all relevant transformed terms
are $\infty$.

If $\delta=2$, set $G=W^\perp$ (rank $d$) and $H=U^\perp$ (rank
$d+2$).  All non-forced transformed terms share the same $p$-coordinate
$\bar p(G)$.  A term indexed by $e$ has $q$-coordinate
$\bar q(L_e)$ for $L_e=(\cl(Ae))^\perp$, a rank $d+1$ flat in
$[G,H]$.  In a rank-two interval, either two non-forced terms have the
same $L_e$ (hence are equal and all other terms are $\infty$), or
there are three distinct points.  In the three-distinct-points case,
use $q\preceq\mu_M$: choose a basis $S$ of $G$, representatives
$\ell_x,\ell_y,\ell_z$ in the three points, and a complement $C$
outside $H$.  The incidence relation for $q\preceq\mu_M$ with lower
set $S$ and upper set $S\ell_x\ell_y\ell_z C$ has precisely the
three finite $\mu_M$-terms $\bar q(L_x),\bar q(L_y),\bar q(L_z)$.

If $\delta=3$, set $G=W^\perp$ (rank $d-1$) and $H=U^\perp$ (rank
$d+2$).  Define
\[
P_x=(\cl(Ayz))^\perp,\quad P_y=(\cl(Axz))^\perp,\quad
P_z=(\cl(Axy))^\perp
\]
and
\[
Q_x=(\cl(Ax))^\perp,\quad Q_y=(\cl(Ay))^\perp,\quad
Q_z=(\cl(Az))^\perp.
\]
Modularity in the rank-three interval gives that $P_x,P_y,P_z$ are
distinct atoms over $G$ and
\[
Q_x=P_y\vee P_z,\quad Q_y=P_x\vee P_z,\quad Q_z=P_x\vee P_y.
\]
Choosing representatives $a_x,a_y,a_z$ over a basis $S$ of $G$, the
original three-term relation for $p\preceq q$ is exactly the
transformed relation.

Applying this elementary lemma along a one-point completion flag
through $\nu$ shows $\nu^{\perp_M}\preceq\mu_M$ for every
$\nu\preceq\mu_M$.  Thus the formula gives an involution on the
underlying projective valuated quotients representing $\Dr(M)$.

\subsection*{Underlying ordinary polar quotient}

For the nested incidence cryptomorphism one also has to track the
underlying ordinary quotient.  If $N$ is an ordinary quotient of $M$ of
rank $k$, the support defined by the polar construction has the standard
rank formula on flats $X$ of $M$:
\[
\rho_{N^{\perp_M}}(X)=\rk_M(X)-k+\rho_N(X^\perp).
\]
This is the usual dual rank formula for matroids on a modular
geometric lattice (the $q$-matroid/polarity formula).  Its bases are
exactly independent $(r-k)$-sets whose polar rank-$k$ flat spans $N$,
matching the support of $\nu^{\perp_M}$.

If $N_p$ is a quotient of $N_q$, then for flats $Y^\perp\subseteq
X^\perp$ the rank increments satisfy
\[
\rho_{N_p}(X^\perp)-\rho_{N_p}(Y^\perp)
\le
\rho_{N_q}(X^\perp)-\rho_{N_q}(Y^\perp).
\]
Substituting in the displayed polar rank formula gives
\[
\rho_{N_q^{\perp_M}}(Y)-\rho_{N_q^{\perp_M}}(X)
\le
\rho_{N_p^{\perp_M}}(Y)-\rho_{N_p^{\perp_M}}(X),
\]
which is the ordinary strong-map criterion for
$N_q^{\perp_M}$ to be a quotient of $N_p^{\perp_M}$.

\subsection*{Circuit-minimum lemma for relative quotients}

If $w\preceq\mu_M$ has rank $k$, $A$ is independent of size $k-1$,
and $C$ is a circuit of the contraction $M/A$ with $|C|\ge2$, then
\[
\min_{i\in C} w(Ai)
\]
is attained at least twice.  Proof: for each $i\in C$,
$A\cup(C-i)$ is independent and spans the same flat.  Choose a
common complement $D$ such that $A\cup(C-i)\cup D$ is an $M$-basis.
Apply $\mathcal I_{k,r}(w,\mu_M;A,A\cup C\cup D)$.  The finite
$\mu_M$-terms are exactly those indexed by $C$; deleting an element
of $D$ leaves the circuit $C$ over $A$.

\subsection*{Direct full-incidence duality}

Let $p\preceq q\preceq\mu_M$ with ranks $a\le b$.  Having already
proved that transforms are quotients of $\mu_M$, and using the ordinary
polar quotient formula for the underlying matroids, it remains to prove
the nested incidence relations for $q^{\perp_M}\preceq p^{\perp_M}$.
If $b=r$, the lower
transformed rank is zero and there are no incidence relations.  Assume
$b<r$.

Fix $A\subseteq T$ with
\[
|A|=r-b-1,\qquad |T|=r-a+1,
\]
and put $B=T\setminus A$, $m=|B|=b-a+2$, $U=\cl(A)$, $W=\cl(T)$.
If $A$ is dependent all transformed terms are $\infty$.  Otherwise
$\rk U=r-b-1$.  A transformed term can be finite only if $T-i$ is
independent, which requires
$\rk W-\rk U\ge b-a+1=m-1$.  Thus the only nontrivial cases are
$\delta=\rk W-\rk U=m$ and $\delta=m-1$.

Nondegenerate case $\delta=m$: $T$ is independent (this case is empty
if $a=0$).  Let $G=W^\perp$, $H=U^\perp$, of ranks $a-1$ and
$b+1$.  For $i\in B$ set
\[
X_i=(\cl(T-i))^\perp,\qquad Y_i=(\cl(Ai))^\perp.
\]
For $J\subseteq B$, the flats $\cl(A\cup J)$ have ranks
$\rk U+|J|$ and form the Boolean sublattice generated by the frame.
After anti-isomorphism, the $X_i$ are atoms over $G$ and
\[
Y_i=\bigvee_{j\ne i}X_j .
\]
Choose a basis $S$ of $G$ and representatives $x_i\in X_i\setminus G$.
Then $S\cup\{x_i:i\in B\}$ has size $b+1$ and the original incidence
relation $\mathcal I_{a,b}(p,q;S,S\cup\{x_i\})$ has terms
\[
p(Sx_i)+q(S\cup\{x_j:j\ne i\})
=\bar p(X_i)+\bar q(Y_i)
=p^{\perp_M}(T-i)+q^{\perp_M}(Ai).
\]
This is exactly the transformed relation.

Degenerate case $\delta=m-1$: in $M/A$, the set $B$ has nullity one
and hence a unique circuit $C$.  If $C$ is a singleton loop, then all
transformed terms are $\infty$.  If $|C|\ge2$, the indices outside
$C$ leave $T-i$ dependent.  For $i\in C$, $T-i$ is independent with
closure $W$, so the $p^{\perp_M}$-coordinate is the common value
$\bar p(W^\perp)$.  The circuit-minimum lemma applied to the quotient
$q^{\perp_M}\preceq\mu_M$ and to the circuit $C$ of $M/A$ gives that
\[
\min_{i\in C} q^{\perp_M}(Ai)
\]
is attained at least twice.  Adding the common $p$-coordinate proves
the transformed incidence relation.

Thus $p\preceq q$ implies $q^{\perp_M}\preceq p^{\perp_M}$ for all
rank gaps.  This is the full order reversal.

\subsection*{Embedded flats}

For an embedded flat $F$ of rank $f$, $q_F=M/F\oplus U_{0,F}$ has
finite rank $r-f$ flat coordinate at $Z$ iff $Z\vee F=\hat1$.  For an
independent $f$-set $A$ with $Y=\cl(A)$,
\[
q_F^{\perp_M}(A)<\infty
\Longleftrightarrow Y^\perp\vee F=\hat1
\Longleftrightarrow Y^\perp\wedge F=\hat0
\Longleftrightarrow Y\vee F^\perp=\hat1,
\]
using modularity and complementary ranks.  This is exactly the basis
condition for $q_{F^\perp}=M/F^\perp\oplus U_{0,F^\perp}$; all finite
values are zero.  Hence the opposition transform extends the given
flat involution.

\subsection*{Literature and computational notes}

Relevant references: Dress--Wenzel for valuated matroids; Speyer for
tropical linear spaces; Brandt--Eur--Zhang for quotient/incidence and
flag Dressians; Murota for valuated truncation/elongation on
independent sets and $M^\natural$-convexity; Jarra--Lorscheid for
flag matroids over the tropical hyperfield/perfect tracts.  Hirai's
valuated matroids on semimodular lattices are conceptually close to
the modular-lattice duality here.

Computational sanity checks from earlier rounds: full subset-indexed
Plucker and incidence relations were tested for Boolean $U_{4,4}$,
$U_{2,4}$, $U_{2,5}$, Fano $PG(2,2)$ with dot-product polarity,
$U_{2,3}\oplus U_{2,3}$, and random samples in $PG(2,3)$; no
counterexample to flat-coordinate opposition was found.  A Z3 search
in $PG(3,2)$ for rank gap $1$ vs.\ $3$ over values $\{0,1\}$ was
UNSAT.  These computations are supportive only; the proof above is
the direct incidence-duality argument.

\end{document}
